\providecommand{\Z}{\mathbb Z}
\providecommand{\C}{\mathbb C}
\providecommand{\La}{\Lambda}
\providecommand{\Sp}{\operatorname{Sp}}
\providecommand{\Orth}{\operatorname{O}}
\providecommand{\Aut}{\operatorname{Aut}}
\providecommand{\Ran}{\operatorname{Ran}}
\providecommand{\CoHA}{\operatorname{CoHA}}
\providecommand{\Hall}{\mathcal H}
\providecommand{\Prim}{\operatorname{Prim}}
\providecommand{\sdim}{\operatorname{sdim}}
\providecommand{\sdet}{\operatorname{sdet}}
\providecommand{\RHom}{\operatorname{RHom}}
\providecommand{\nuDelta}{\nu_{\Delta_5}}
\providecommand{\epsor}{\epsilon_o}
\providecommand{\gDelta}{\mathfrak g_{\Delta_5}}
\providecommand{\WII}{W^{(2)}(\La^{2,1}_{II})}
\providecommand{\KX}{\mathfrak K_X}
\providecommand{\FX}{\mathcal F_X}
\providecommand{\HX}{\mathcal H_X}

\section{Agent 15-E: the character comparison
\texorpdfstring{\(\epsilon_o=\nu_{\Delta_5}\)}{epsilon equals nu Delta5}}
\label{sec:agent15-epsilon-nu-delta5}

\subsection{Attacked equality}

The attacked statement is the strongest possible gate-5 assertion:
\[
        \epsor=\nuDelta .
\]
In this form it is not a theorem and not even a typed assertion.  The
right side is the Maass automorphic character of the Siegel modular form
\(\Delta_5\).  It is already defined on \(\Sp_4(\Z)\), and its
restriction to the type-II Weyl group is theorem-level.  The left side
is supposed to be the monodromy character of a chosen square root of the
virtual determinant line in an oriented compact \(K3\times E\)
Hall/CoHA theory.  That character is not defined until the compact
source, hybrid \(\Ran(E)\) object, oriented Hall atlas, protected
integration, and Weyl lift have been constructed.

The false shortcut is:
\[
\hbox{the scalar determinant is } \Delta_5
\quad\Longrightarrow\quad
\hbox{the compact orientation character is } \nuDelta .
\]
This implication fails.  The scalar determinant supplies a modular form
and a product.  It does not supply a determinant-line square root on the
compact Hall stack, a loop in the oriented Hall moduli problem, or the
transport sign of that square root around the loop.

\subsection{Status after attacking all six gates}

\begin{center}
\small
\begin{tabular}{p{0.18\textwidth}p{0.34\textwidth}p{0.38\textwidth}}
\hline
Gate & attack & healed status \\
\hline
1. compact \(E_3\) source &
No constructed compact finite-first holomorphic \(E_3\) source on
\(K3\times E\) is present in the inspected material. &
The source must be supplied with charge completion, QME/anomaly control,
descent, and the Lorentzian degree map
\(\alpha(n,l,m)=2nf_2-lf_3+2mf_{-2}\).  Without this, there is no
source-side action on charges to compare with a Weyl reflection. \\
2. hybrid \(\Ran(E)\) wrapping &
Ordinary finite-support Ran locality misses positive elliptic degree:
a curve of positive elliptic degree projects onto all of \(E\). &
Gate 5 must be computed on the eventual hybrid local/global wrapped
object, not on a naive support-local Ran algebra.  Mixed local-wrapped
extension stacks are part of the sign problem. \\
3. oriented compact Hall/CoHA atlas &
The orientation character is undefined without a strong orientation
system and oriented extension correspondences. &
One needs square roots of determinant lines, vanishing cycles,
Thom--Sebastiani transport, finite HN completions, and a lifted
\(\WII\)-action on the oriented completed Hall object. \\
4. protected integration &
A scalar protected trace can forget the square-root monodromy. &
The protected integration map must be equivariant for the lifted Weyl
action and must preserve the orientation branch before taking the scalar
shadow. \\
5. character comparison &
The equality is not produced by normalizing \(\Delta_5\), by
\(\Delta_5^2\), or by OP/DT constants. &
The correct theorem is conditional: after gates 1--4 supply the
Hall-side domain and the three type-II reflection signs are computed,
\(\epsor=\nuDelta|_{\WII}\). \\
6. primitive recognition &
The equality of characters does not identify a BKM bracket. &
If gate 5 closes, it supplies the reflection-character compatibility
needed by the primitive recognition theorem.  The bracket, parity split,
simple primitive representatives, Hopf pairing, and radical quotient
remain separate hypotheses. \\
\hline
\end{tabular}
\end{center}

\subsection{Domain problem}

The automorphic character is
\[
        \nuDelta:\Sp_4(\Z)\longrightarrow \{\pm1\},
        \qquad
        \Delta_5|_5 g=\nuDelta(g)\Delta_5 .
\]
In the manuscript it is given by Maass's formulas on the standard
generators:
\[
\nuDelta(J)=1,\qquad
\nuDelta(T(B))=(-1)^{b_1+b_2+b_3},
\]
\[
\nuDelta(g(A))
=(-1)^{(1+a_1+a_4)(1+a_2+a_3)+a_1a_4}.
\]
The central element \(-I\) has Maass value \(+1\) in this normalization,
so the character descends to \(\Sp_4(\Z)/\{\pm I\}\) for the present
orthogonal comparison.

The Hall-side character would have the form
\[
        \epsor:\WII\longrightarrow \{\pm1\}.
\]
It is not a scalar extracted from the partition function.  It is the
sign of orientation-line transport.  A type-correct definition requires
the following typed package:
\[
(\HX^{\mathrm{or}},o,I^{\mathrm{prot}},\rho_o),
\]
where \(o\) is a strong square-root orientation system, \(I^{\mathrm{prot}}\)
is protected integration compatible with the Hall product and HN
completion, and
\[
        \rho_o:\WII\longrightarrow
        \Aut(\HX^{\mathrm{or}},o,I^{\mathrm{prot}})
\]
is a lifted Weyl action preserving the charge-to-root map.  Then
\[
        \rho_o(w)^*\ell_{\gamma}
        \simeq
        \epsor(w)\,\ell_{w\gamma}
\]
defines \(\epsor(w)\).  Without \(\rho_o\) the symbol \(\epsor(w)\) has
no argument.  Without \(o\) it has no line on which to act.  Without
protected integration it is not tied to the protected BPS trace.  Without
the charge map it is not the same reflection seen by the automorphic
form.

\subsection{Exact comparison diagram}

The comparison is an equality of two characters on one group after both
vertical maps are constructed:
\[
\begin{array}{ccc}
\WII
& \xrightarrow{\ \iota_{\mathrm{ext}}\ } &
\Sp_4(\Z)/\{\pm I\}
\\[0.7em]
\rho_o \downarrow
&&
\downarrow \overline{\nuDelta}
\\[0.7em]
\Aut(\HX^{\mathrm{or}},o,I^{\mathrm{prot}})
    /\Aut^+(\HX^{\mathrm{or}},o,I^{\mathrm{prot}})
& \xrightarrow{\ \operatorname{sgn}_o\ } &
\{\pm1\}.
\end{array}
\]
Here \(\iota_{\mathrm{ext}}\) is the exterior-square/orthogonal
identification used by the manuscript:
\[
\WII\subset \Orth(\La^{2,1})_+
\hookrightarrow \Orth(\La^{3,2})_+
\simeq \Sp_4(\Z)/\{\pm I\}.
\]
The bottom character is
\[
        \epsor=\operatorname{sgn}_o\circ\rho_o.
\]
The target equality is therefore
\[
        \operatorname{sgn}_o\circ\rho_o
        =
        \overline{\nuDelta}\circ\iota_{\mathrm{ext}}
        \qquad\hbox{on }\WII .
\]
This diagram is the minimum type-correct form of gate 5.

\subsection{Automorphic side: theorem-level computation}

The type-II chamber has simple real walls
\[
        \delta_1=2f_2-f_3,\qquad
        \delta_2=2f_{-2}-f_3,\qquad
        \delta_3=f_3 .
\]
The manuscript computes
\[
        \nuDelta(s_{\delta_i})=-1,\qquad i=1,2,3,
\]
and
\[
        \nuDelta(s_{f_{-2}-f_2})
        =
        \nuDelta(s_{f_2-f_3})
        =
        \nuDelta(s_{f_2+f_3})=1 .
\]
Consequently, in the type-II Levi decomposition
\[
        \Orth(\La^{2,1})_+
        \simeq
        \WII\rtimes\Aut(\Poly_{II}),
        \qquad \Aut(\Poly_{II})\simeq S_3,
\]
one has
\[
        \Delta_5(w\cdot a Z)=\det(w)\Delta_5(Z),
        \qquad
        w\in \WII,\ a\in\Aut(\Poly_{II}).
\]
Thus the automorphic side of the desired equality is:
\[
        \overline{\nuDelta}\circ\iota_{\mathrm{ext}}
        =
        \det|_{\WII}.
\]
This is a theorem-level statement from Maass's character, Nikulin's
reflection group, and the exterior-square identification.  It is not a
Hall orientation statement.

\subsection{Generator-level tests for the Hall side}

The compact orientation comparison must pass the following tests before
gate 5 can be promoted from conjecture to theorem.

\begin{enumerate}
\item \textbf{Domain test.}  Construct
\[
\rho_o:\WII\to\Aut(\HX^{\mathrm{or}},o,I^{\mathrm{prot}})
\]
on the completed Hall object, not merely on its scalar partition
function.  The action must preserve the HN completion and the mixed
local/global wrapped sectors.

\item \textbf{Charge test.}  The action must intertwine the compact
charge grading with
\[
        \alpha(n,l,m)=2nf_2-lf_3+2mf_{-2}.
\]
The reflection \(s_{\delta_i}\) on the Hall side must be the same
reflection that \(\iota_{\mathrm{ext}}\) sends to the Maass generator on
the automorphic side.

\item \textbf{Three real-wall tests.}  For each type-II real wall
\[
\delta_1,\delta_2,\delta_3,
\]
transport the chosen determinant-line square root around the corresponding
simple diagonal wall loop and prove
\[
        \epsor(s_{\delta_1})
        =
        \epsor(s_{\delta_2})
        =
        \epsor(s_{\delta_3})
        =
        -1 .
\]
The expected local model is Pfaffian: the simple real-root reflection
reverses one oriented normal mode.  This is a geometric lemma, not a
theta-normalization convention.

\item \textbf{Involution test.}  The transported orientation around
\(s_{\delta_i}^2\) must be \(+1\).  A sign \(-1\) on the square would
show that \(\rho_o\) is only a projective action or that the orientation
line has been transported through the wrong cover.

\item \textbf{Relation test.}  The computed signs must respect the
relations in Nikulin's reflection group.  Equivalently, the orientation
transport must be an honest character
\[
        \epsor:\WII\to\{\pm1\},
\]
not a \(2\)-cocycle on a projective Weyl action.

\item \textbf{\(S_3\)-chamber automorphism test.}  The chamber
automorphisms generated by \(s_{f_2-f_3}\) and \(s_{f_{-2}-f_2}\) have
automorphic sign \(+1\).  The Hall orientation must not introduce a
spurious sign when these elements merely permute the three null vertices
of \(\Poly_{II}\).

\item \textbf{Type-I control test.}  Do not test the equality on the
wrong generators.  In the type-I decomposition the sign is carried by
\(\Aut(\Poly_I)\), while type-I reflections lie in the Maass kernel.
Confusing type I and type II gives the wrong character.

\item \textbf{OP chamber/source normalization test.}  The OP scalar
branch uses
\[
        D_5=64^{-1}\Delta_5,\qquad
        Z^{K3\times E}_{\mathrm{OP}}
        =
        -4096\,\Delta_5^{-2}.
\]
These constants must not enter the orientation computation.  They fix
the scalar branch and the theta-to-monic normalization, not the
Hall-side square-root monodromy.

\item \textbf{Wrapped-sector test.}  For positive elliptic degree, the
object dominates \(E\).  The sign must therefore be computed after the
hybrid local/global wrapped extension correspondence exists.  A
projection-finite Ran test alone cannot decide the full Igusa
\(s\)-direction.

\item \textbf{Protected integration equivariance test.}  The protected
integration map must satisfy
\[
        I^{\mathrm{prot}}(\rho_o(w)x)
        =
        \epsor(w)\,w\cdot I^{\mathrm{prot}}(x)
\]
on the oriented completed Hall object.  A scalar equality after
forgetting orientation is too weak.
\end{enumerate}

\subsection{Sign and normalization failure modes}

\begin{enumerate}
\item \textbf{Scalar-rescaling failure.}
Multiplying \(\Delta_5\) by \(64^{-1}\), or by any nonzero scalar,
does not change \(\nuDelta\).  It cannot construct \(\epsor\).

\item \textbf{Square failure.}
\(\Delta_5^2=\Delta_{10}\) is scalar.  The square forgets the sign
being tested.  The OP/DT identity
\[
        Z_{\mathrm{OP}}^{K3\times E}
        =
        Z_{\mathrm{DT}}^{K3\times E}
        =
        -4096\,\Delta_5^{-2}
\]
is a scalar branch, not an orientation monodromy computation.

\item \textbf{Behrend-sign failure.}
The Behrend-weighted Euler characteristic gives a protected scalar
shadow.  It does not by itself produce the square-root orientation line
or its Weyl monodromy.

\item \textbf{Exterior-square lift failure.}
The automorphic side lives on \(\Sp_4(\Z)\) and descends through
\(\{\pm I\}\) in the Maass normalization.  A Hall-side Weyl action must
choose the same exterior-square lift.  A hidden central sign would change
the comparison diagram.

\item \textbf{Orientation-twist failure.}
A strong orientation is a torsor.  Even after orientability is known,
an equivariant orientation can be twisted by a character unless a base
branch and its Weyl transport are fixed.  Therefore orientability of
\(\operatorname{RPerf}(K3\times E)\) is not the equality
\(\epsor=\nuDelta\).

\item \textbf{Quotient-first failure.}
Reducing by the elliptic translation before Hall correspondences are
formed can erase the relative \(E\)-position needed to define mixed
extension stacks and their determinant-line transports.

\item \textbf{Analogy failure.}
The theta-characteristic explanation of the Maass sign is a guide to the
answer.  It is not a proof of determinant-line monodromy for the compact
Hall orientation.
\end{enumerate}

\subsection{Strongest healed theorem}

\begin{theorem}[conditional character comparison]
Let \(X=K3\times E\).  Assume the compact Igusa package supplies:
\begin{enumerate}
\item a finite-first charge-completed holomorphic \(E_3\)-source
\(\mathcal A_X^{E_3}\) with anomaly/QME/descent control;
\item a chain-level \(K3\)-to-\(E\) specialization to a hybrid
local/global wrapped factorization object \(\FX\);
\item an oriented compact critical Hall/CoHA object \(\HX^{\mathrm{or}}\)
with vanishing cycles, extension correspondences, HN completion, and
Thom--Sebastiani orientation transport;
\item a protected integration map \(I^{\mathrm{prot}}\) compatible with
the completed Hall product and wall crossing;
\item a lifted Weyl action
\[
        \rho_o:\WII\to
        \Aut(\HX^{\mathrm{or}},o,I^{\mathrm{prot}})
\]
preserving the charge map \(\alpha(n,l,m)=2nf_2-lf_3+2mf_{-2}\);
\item exterior-square compatibility
\[
        \iota_{\mathrm{ext}}\circ\rho_o
\]
with the automorphic reflection action used to compute
\(\nuDelta\);
\item the three Pfaffian normal-mode computations
\[
        \operatorname{sgn}_o(\rho_o(s_{\delta_i}))=-1,
        \qquad i=1,2,3,
\]
and the relation/cocycle check showing that these signs define an honest
character of \(\WII\).
\end{enumerate}
Then
\[
        \epsor
        =
        \nuDelta|_{\WII}.
\]
\end{theorem}

\begin{proof}
By Maass's formula and the exterior-square/lattice identification,
\(\nuDelta|_{\WII}=\det|_{\WII}\), with value \(-1\) on the three
simple type-II reflections.  By the hypotheses, the Hall-side
orientation transport gives an honest character \(\epsor\) of the same
group, for the same reflections, with the same values on the three
simple type-II reflections.  Since the lifted action satisfies the
reflection-group relations, the two characters agree on the group
generated by these reflections, namely \(\WII\).
\end{proof}

This is the strongest proved-form statement available after the attack.
Its hypotheses are not all constructed in the current paper or in the
surrounding volumes.

\subsection{Healed conjecture for the compact programme}

\begin{conjecture}[compact orientation monodromy]
For the actual compact \(K3\times E\) Hall/CoHA realization, once gates
1--4 have been constructed, the Pfaffian normal-mode computation holds:
transporting the strong orientation around each simple type-II diagonal
wall reflection reverses exactly one oriented normal factor.  Hence
\[
        \epsor(s_{\delta_i})=-1,\qquad i=1,2,3,
\]
and therefore
\[
        \epsor=\nuDelta|_{\WII}.
\]
\end{conjecture}

This conjecture is sharp.  It names the missing geometric computation.
It does not infer the orientation character from the scalar determinant.

\subsection{Dependence on gates 1--4}

Gate 5 depends on gates 1--4 in a strict order.

\begin{enumerate}
\item Gate 1 supplies the compact charge-completed \(E_3\) source and
the charge lattice.  Without it there is no compact object on which
reflections can act.
\item Gate 2 supplies the hybrid wrapped factorization object over
\(E\).  Without it the positive elliptic-degree part of the Igusa
product has no locality domain.
\item Gate 3 supplies the orientation lines and extension
correspondences.  Without it \(\epsor\) is undefined.
\item Gate 4 supplies protected integration and equivariance.  Without
it the sign is not tied to the protected BPS trace or to wall crossing.
\end{enumerate}

Thus gate 5 cannot be made theorem-level by strengthening only the
automorphic section of the paper.  The missing work is geometric.

\subsection{Consequence for gate 6}

If gate 5 is proved, it gives the reflection-character compatibility
needed by primitive recognition:
\[
        \epsor|_{\WII}=\nuDelta|_{\WII}.
\]
This is necessary for the compact primitive denominator to transform as
the Gritsenko--Nikulin denominator.  It is not sufficient for
\[
        H^0\Prim_{\mathrm{prot}}(\FX)/\operatorname{Rad}
        \cong \gDelta .
\]
Gate 6 still needs:
\begin{enumerate}
\item the compact protected primitive functor;
\item full parity dimensions, not only signed superdimensions;
\item simple primitive representatives for the real and imaginary simple
roots;
\item the Chevalley, real Serre, Borcherds orthogonality, and super sign
relations;
\item generation by simple primitives and the Cartan;
\item a completed Hopf pairing with the same radical quotient as the
Gritsenko--Nikulin/Kac construction.
\end{enumerate}
The character comparison can certify the denominator anti-invariance.
It cannot manufacture the Hall bracket.

\subsection{Local and cross-volume anchors}

\begin{itemize}
\item \verb|CLAUDE.md:1--39| and \verb|AGENTS.md:1--118| fix the
repository role, no-destructive-git rule, proof-grade claim discipline,
and cross-volume consistency requirement.

\item \verb|main.tex:656--686| defines \(\Delta_5\), the Maass character
\(\nuDelta\), and the scalar square \(\Delta_5^2=\Delta_{10}\).

\item \verb|main.tex:722--735| records the diagonal divisor and names
\(\epsilon_{\det}:=\nuDelta\) as the automorphic determinant sign.

\item \verb|main.tex:906--933| states the protected Hall integration
criterion and requires the orientation character to match \(\nuDelta\) on
the type-II Weyl group.

\item \verb|main.tex:936--955| states the orientation-monodromy open
problem and explicitly says the character is not produced by normalizing
\(\Delta_5\).

\item \verb|main.tex:1436--1458| computes the Maass signs on the simple
reflection classes.

\item \verb|main.tex:1485--1545| records Nikulin's reflection group,
the type-I/type-II chamber decomposition, and the multiplier restriction.

\item \verb|main.tex:2258--2342| defines the compact Igusa realization
datum, including the \(E_3\) source, specialization, oriented Hall
object, chiral coalgebra, Koszul comparison, and primitive recognition
certificate.

\item \verb|main.tex:2935--2998| separates the denominator theorem, the
OP/DT scalar branch, and the compact realization consequences.

\item \verb|proj.bib:15--21| is Maass 1964 for the Siegel multiplier.
\verb|proj.bib:64--98| are the Gritsenko--Nikulin automorphic
correction references.  \verb|proj.bib:253--274| are the
Kontsevich--Soibelman and Davison--Meinhardt CoHA references.

\item
\verb|agent_material/14_orientations_protected_integration_from_volumes_attack_heal.tex:21--148|
is the immediate predecessor report: it distinguishes the automorphic
character from the undefined geometric orientation character and gives
the type-correct diagram.

\item
\verb|agent_material/14_ranE_hybrid_wrapping_from_volumes_attack_heal.tex:23--88|
is the locality obstruction for positive elliptic degree.

\item
\verb|agent_material/14_compact_E3_source_from_volumes_attack_heal.tex:16--33|
states that no compact \(E_3\) source theorem is present in the volumes.

\item
\verb|agent_material/14_primitive_borcherds_recognition_from_volumes_attack_heal.tex:18--37|
and
\verb|agent_material/14_primitive_borcherds_recognition_from_volumes_attack_heal.tex:140--246|
separate primitive recognition from signed determinant data.

\item
\verb|agent_material/14_chiral_bar_cobar_Koszul_realization_attack_heal.tex:23--95|
states that bar--cobar inversion starts after the chiral object has been
supplied.

\item
\verb|/Users/raeez/calabi-yau-quantum-groups/FRONTIER.md:27--31|
states the two-stage \(\Phi_d=\mathrm{Sp}\circ\Phi_d^{FA}\)
architecture; \verb|FRONTIER.md:81--88| records bracket-level
BPS--BKM comparison as open.

\item
\verb|/Users/raeez/calabi-yau-quantum-groups/chapters/examples/cy_d_kappa_stratification.tex:3750--3784|
checks the Maass multiplier through Steinberg, Weil, and orthogonal
routes; \verb|:4061--4090| records the type-I/type-II sign asymmetry.

\item
\verb|/Users/raeez/calabi-yau-quantum-groups/chapters/theory/cy3_chain_level_bridge.tex:522--701|
defines the oriented hCS--Hall comparison datum and says the theorem
starts only after orientation, trace, twist, completion, product, and
coherence data exist.

\item
\verb|/Users/raeez/calabi-yau-quantum-groups/chapters/theory/cy3_chain_level_bridge.tex:1099--1211|
defines the hCS--Hall descent obstruction, including the determinant-line
square-root mismatch \(o_{\mathrm{or}}\).

\item
\verb|/Users/raeez/calabi-yau-quantum-groups/chapters/theory/cy3_chain_level_bridge.tex:3452--3695|
proves Hall-side orientation torsor triviality and a finite DWR/Ran
comparison only after finite primitive packages are supplied; the
pro-limit obstruction remains.

\item
\verb|/Users/raeez/calabi-yau-quantum-groups/chapters/examples/k3e_bkm_chapter.tex:14082--14106|
separates theorem-grade numerical \(-\Delta_5^{-2}\) from the conditional
compact critical CoHA reading.

\item
\verb|/Users/raeez/chiral-bar-cobar/CLAUDE.md:56--63| says Vol. I's
bar--cobar theorem is the \(E_1\)-chiral shadow after Stage-2
specialization, not the compact \(E_3\) source.

\item
\verb|/Users/raeez/chiral-bar-cobar-vol2/chapters/theory/factorization_swiss_cheese.tex:35--56|
states the local operad is the formal-disc shadow of global
factorization data, not an inverse construction.

\item
\verb|/Users/raeez/chiral-bar-cobar-vol2/chapters/connections/ht_bulk_boundary_line_core.tex:1030--1049|
proves a boundary Thom--Sebastiani multiplicativity theorem; it is not a
compact \(K3\times E\) orientation-monodromy computation.

\item
\verb|/Users/raeez/topological-strings/main.tex:2434--2479|
imports Costello/Costello--Li BV technology but records the
specialization problem; \verb|:2624--2646| separates algebraic product
targets from analytic graph construction.
\end{itemize}

\subsection{Files changed}

Only this file was written:
\verb|agent_material/15_epsilon_equals_nu_Delta5_character_comparison_attack_heal.tex|.

\subsection{Checks run}

Checks run:
\begin{verbatim}
git diff --no-index --check -- /dev/null agent_material/15_epsilon_equals_nu_Delta5_character_comparison_attack_heal.tex
git diff --cached --check -- agent_material/15_epsilon_equals_nu_Delta5_character_comparison_attack_heal.tex
rg -n "(Attacked equality|Domain problem|Exact comparison diagram|Generator-level tests|Sign and normalization failure modes|Strongest healed theorem|Dependence on gates 1--4|Consequence for gate 6|Local and cross-volume anchors|Files changed|Checks run|Open questions)" agent_material/15_epsilon_equals_nu_Delta5_character_comparison_attack_heal.tex
\end{verbatim}
The whitespace checks reported no diagnostics, and the heading scan found
all required report sections.  No full
manuscript build is required for this single agent report unless the
main thread integrates it into \texttt{main.tex}.

\subsection{Open questions}

\begin{enumerate}
\item What is the exact compact moduli problem whose determinant line
will carry the orientation?  Stable pairs, sheaves, objects of
\(\operatorname{Perf}(K3\times E)\), or a hybrid Hall heart give
different determinant-line loops.

\item Does the type-II Weyl group act by actual autoequivalences or
correspondences on the compact oriented Hall object, or only on the
automorphic/root lattice shadow?

\item Which loop in the compact Hall stack represents the simple
reflection \(s_{\delta_i}\)?  The answer must be stated before the
Pfaffian normal-mode sign can be computed.

\item Does the determinant-line transport around \(s_{\delta_i}\) reverse
exactly one normal mode, or is there an additional contribution from
wrapped elliptic-degree sectors, the \(E\)-quotient, or the HN
completion?

\item Does the orientation transport satisfy the Nikulin reflection-group
relations strictly, or only projectively?  A projective defect would
replace \(\epsor\) by a \(2\)-cocycle and break the comparison.

\item Can the finite DWR/Ran oriented comparison packages in Vol. III be
made compatible under the \(N,r,L,m\) inverse systems, killing the
displayed \(\varprojlim^1\) obstruction?

\item Once gate 5 is closed, can gate 6 exhibit the actual primitive
generators and Borcherds--Kac relations, rather than only the
denominator anti-invariance?
\end{enumerate}
